\section{Nghiên cứu người dùng}

\subsection{Đối tượng người dùng mục tiêu}
% Trình bày:
% - Người dùng mục tiêu (Chủ nuôi thú cưng bận rộn theo Proposal)
% - Đặc điểm nhân khẩu học và thói quen chăm sóc thú cưng
% - Ngữ cảnh sử dụng thực tế (Context of use)

\subsection{Phương pháp nghiên cứu}
% Trình bày phương pháp thực hiện (Phỏng vấn sâu, Khảo sát, Quan sát...):
% - Mục đích
% - Số lượng và tiêu chí người tham gia
% - Quy trình thực hiện
% - Dữ liệu thu thập

\subsection{Kết quả nghiên cứu}
% Trình bày các phát hiện và dữ liệu thực tế:
% - Người dùng thường gặp khó khăn khi...
% - Người dùng kỳ vọng...
% - Người dùng nhầm lẫn hoặc lo lắng về...

\subsection{Nhu cầu và điểm đau của người dùng}
% Tổng hợp nhu cầu và điểm đau chính của người dùng (kế thừa từ Proposal):
% - Điểm đau: Chờ xác nhận lâu, thất lạc dặn dò đặc biệt/dị ứng, thiếu minh bạch tiến độ...
% - Nhu cầu: Đặt lịch tức thì, tự động lưu hồ sơ dị ứng/thuốc, theo dõi tiến độ thời gian thực...

\subsection{Chân dung người dùng}
% Bắt buộc trình bày ĐẦY ĐỦ TẤT CẢ các Persona có trong dự án (không được bỏ sót bất kỳ Persona nào):
% - Persona 1: Nguyễn Hoàng Lan (26 tuổi, Chuyên viên Marketing) - Chủ nuôi cẩn trọng, nuôi mèo Mochi nhút nhát có tiền sử dị ứng sữa tắm
% - Persona 2: Lê Hoàng Minh (22 tuổi, Sinh viên / Lập trình viên) - Chủ nuôi bận rộn, nuôi 2 thú cưng (chó Poodle Bơ & mèo Miu)
% Với mỗi Persona, trình bày đầy đủ: Bối cảnh, Mục tiêu (Goals), Hành vi (Behaviors), Điểm đau (Pain points), Nhu cầu (Needs).

\subsection{Đề xuất giá trị}
% Bắt buộc trình bày Đề xuất giá trị (Value Proposition Canvas) ĐỐI ỨNG 1-1 TƯƠNG ỨNG VỚI TỪNG PERSONA:
% (Nguyên tắc: Có bao nhiêu Persona thì phải có bấy nhiêu bản Đề xuất giá trị tương ứng)
% - Đề xuất giá trị cho Persona 1 (Nguyễn Hoàng Lan):
%   + Hồ sơ khách hàng (Customer Profile): Tác vụ (Jobs), Điểm đau (Pains), Điểm mong muốn (Gains)
%   + Bản đồ giá trị (Value Map): Sản phẩm & Dịch vụ (Products & Services), Trợ thủ giải tỏa đau đớn (Pain Relievers), Yếu tố tạo lợi ích (Gain Creators)
% - Đề xuất giá trị cho Persona 2 (Lê Hoàng Minh):
%   + Hồ sơ khách hàng (Customer Profile): Tác vụ (Jobs), Điểm đau (Pains), Điểm mong muốn (Gains)
%   + Bản đồ giá trị (Value Map): Sản phẩm & Dịch vụ (Products & Services), Trợ thủ giải tỏa đau đớn (Pain Relievers), Yếu tố tạo lợi ích (Gain Creators)

